\scanpage{225}
\margin{\deleted{1 \unsure{cas initial} $X$ \unsure{du type} \illegible{2} \unsure{corps de base}}}
On a vu plus haut qu’une place est \unsure{trs|très} loin d’être caractérisée par le
foncteur fibre correspondant, \deleted{\unsure{le}}
considéré comme foncteur dans les $\mathbf Z$-modules.
Mais que se passe-t-il si on le regarde comme foncteur dans les structures de Hodge ?
\deleted{\unsure{Ceci devrait être} \unsure{disons que}}
\inserted{Vrai par un théorème \unsure{Rauch}, et par \illegible{1}}

\deleted{\unsure{par ailleurs le point}, on est \unsure{ramené} au cas où
\unsure{deux} \illegible{4} sur $k_0=$ clôture alg. de $\mathbf Q$ dans $k$,
et \unsure{puisqu’on a} $X$ est de type
$X$ est de type fini sur \deleted{\unsure{un}} $\mathbf Q$,
puis par changement de base on \unsure{a} $X$ est connexe de type fini sur $\mathbf C$,
et qu’on a deux pts \inserted{ordinaires} de $X$ à val. dans $\mathbf C$.
Supposons pour simplifier que \deleted{\unsure{tt}} $X$ est
\deleted{affine}, \inserted{(quasi-projectif)},
\unsure{voire affine}, \unsure{commençons par un cas simple} $X=\mathbf P^1_{\mathbf C}$.}

\margin{Attention, il faut partir de pts complexes, \unsure{non seulement de}
places, pour pouvoir définir la structure de Hodge.
\diagram{\begin{tikzcd}[column sep=small,row sep=small]
 & X \arrow[dl,no head] \arrow[dr,no head] & \\
 \unsure{\mathbf Q} \arrow[rr,no head] && k
\end{tikzcd}}
\deleted{$\unsure{A_i}$}
\figure{\begin{tikzpicture}\node at (0,.8) {$X$};\draw (0,.55)--(0,0);\end{tikzpicture}}}

Il semble qu’on le vérifiera facilement dans le cas de
$X=\mathbf P^1_{\mathbf C}-(0)-(1)-(\infty)$,
\inserted{(et deux pts à val. dans $\mathbf C$ \unsure{ord.})},
où les théorèmes classiques \deleted{\unsure{dit}} le donnent ---
[prendre comme motif celui défini par un schéma en courbes elliptiques sur $X$,
d’invariant l’indéterminée $j$].
\margin{224}
